Hoewel er steeds meer aan het licht komt, toont een reality check dat mijn brein vooral in het duister tast

Het brein wil het leven uitvogelen

dus de ziel deel is nep want de geest is er maar wel een deel dat is verbonden met de ziel. Die verbinding met God is de ziel

Ik krijg morgen niet opgelost laat staan gisteren. Alleen nu is er perfectie

Ik mis het vertrouwen om de toekomst en verleden los te laten. Het is waarom het brein is ontwikkeld om de mens

Rembo \& Rembo, Witte pak met stront, Het witte pak staat voor authenticiteit en de stront van wat anderen op je pak smeren, hun viezigheid die ze zelf niet willen opruimen, metafoor voor sociale of morele bezoedeling (denk aan het bijbelse “de zonden afwassen”), in hedendaagse literatuur betekent authenticiteit vaak kwetsbaarheid, die gemakkelijk door anderen misbruikt of bezoedeld kan worden, – je authentieke zelf (wit pak) wordt zichtbaar en daardoor kwetsbaar
– anderen smeren hun eigen onopgeloste problemen (stront) aan op jouw authenticiteit
– het besmeurde pak symboliseert hoe oprechtheid en openheid risico lopen bezoedeld te worden door de omgeving, een houding uitstralen van laat me met rust, een absurdistisch, over-the-top toneel neer; daardoor communiceren ze indirect: wij trekken ons niks aan van jouw normen, hou je oordeel maar bij jezelf, vergelijkbare walging van bemoeizucht zie je in literatuur en theater van de absurde traditie (beckett, ionesco), waarin het individu gefrustreerd raakt door de eisen en vuiligheid van de buitenwereld, ‘laat me met rust’ is soms een schreeuw om ruimte voor je eigen innerlijke dialoog
– het omarmen van die positie in creatief werk betekent: niet langer verdedigen, maar gewoon bestaan—zelfs met stront op het pak
– het is de erkenning dat je niet alles hoeft op te ruimen wat een ander op jou projecteert, de uitspraak “dat is toch niet normaal” die je in de context van de gluurbuur hoort, voegt een belangrijke laag toe aan het symbolische en psychologische spel dat je schetst.

logische en thematische analyse
– de gluurbuur fungeert niet alleen als observator, maar als spreekbuis voor het collectieve oordeel: telkens wanneer er iets buiten het gangbare gebeurt, volgt die afkeurende reactie
– de makers en personages, met hun “laat me met rust”-houding, zetten zich hiermee af tegen de constante neiging om het abnormale te benoemen, te veroordelen en te marginaliseren, “dat is toch niet normaal” is een projectie van schaamte of angst voor het afwijkende; het fungeert als ontkenning van eigen verlangens of vuiligheid
– de gluurbuur wil niet geconfronteerd worden met zijn eigen onopgeruimde rommel, dus wijst die extra duidelijk naar het ‘abnormale’ gedrag van de ander
– hierdoor ontstaat een dynamiek waarin authenticiteit of opstandigheid dubbel wordt gestraft: niet alleen krijg je stront op je pak, je wordt ook nog eens als zonderling gebrandmerkt,

Terwijl de gluurbuur zelf abnormaal gedrag vertoont en verwijtbaar door te gluren in de hoop mensen naakt te zien uit onbewuste lust.

seks is een van de krachtigste energieën in het menselijk bestaan. het is diep geworteld in het lichaam, maar het is ook door het denken meegenomen naar het domein van verbeelding, angst, schuld en verlangen. we kijken zelden naar seks zoals het werkelijk is; we dragen eeuwen aan conditionering met ons mee, van religieuze aard, van culturele herkomst, of persoonlijk. in veel delen van de wereld wordt seks geassocieerd met zonde, in andere worden het verheerlijkt en als handelswaar verkocht. maar zeer weinigen van ons kijken deze kracht direct aan, zonder oordeel, zonder weerstand, zonder overgave.

ji. krishnamurti werd vaak gevraagd naar seks. mensen kwamen tot hem in verwarring: monniken die gelofte van celibaat hadden afgelegd, jonge mensen verscheurd tussen verlangen en schuld, zoekers die dachten dat ze seksualiteit moesten opgeven om tot waarheid te komen. aan ieder van hen wees krishnamurti niet op een direct antwoord, maar op een dieper zien. hij moraliseerde niet, hij zei nooit: doe dit, of: doe dat niet. in plaats daarvan vroeg hij: “heb je verlangen zelf begrepen?” niet enkel seksueel verlangen, maar verlangen in al zijn vormen.

seks is niet louter een biologische daad, het is niet slechts het voortbestaan van de soort. voor velen wordt het het enige moment van volledige overgave, van jezelf vergeten, al is het maar voor even. krishnamurti merkte op: op dat moment is er geen zelf. dat maakt seks zo belangrijk, niet door het fysieke aspect, maar om de psychologische ontsnapping die het biedt. de geest, die voortdurend kletst en streeft, wordt ineens kort stil; voor een ogenblik is er geen strijd, geen worden, slechts zijn.

daarna komt het denken. het herinnert zich het genot, maakt er een beeld van, en dat beeld wordt het centrum van verlangen. zo wordt verlangen geboren. dus is de vraag niet of seks goed of slecht is – dat is een oppervlakkige vraag. de echte vraag is: wat is de aard van verlangen? kun je het waarnemen, niet als iets dat onderdrukt of geconsumeerd moet worden, maar als iets dat puur geobserveerd wordt? als er honger is, eet je; als er verlangen is, kun je erbij blijven, niet vluchten in fantasie of in schuldgevoel, maar eenvoudig blijven en waarnemen.

volgens krishnamurti is deze waarneming het begin van intelligentie. hij maakte een duidelijk onderscheid tussen celibaat en kuisheid. celibaat is vaak opgelegd, het wordt een discipline, een gelofte, een onderdrukking. en wat onderdrukt wordt, wordt alleen maar sterker; het werkt in het donker, via dromen en onderbewuste verlangens. kuisheid is iets geheel anders: het is de zuiverheid van een geest die de beweging van denken en verlangen heeft begrepen. kuisheid is niet het resultaat van inspanning of controle; het is de geur van inzicht.

religies hebben geprobeerd seks te beheersen met angst en dogma. maar krishnamurti zag dat onderdrukking alleen maar conflict veroorzaakt. je kunt misschien niet handelen naar je verlangen, maar de innerlijke storm blijft doorgaan. zolang je met jezelf vecht, kan er geen vrede of vrijheid zijn. hij zei ooit: om vrij van iets te zijn, moet je het begrijpen, niet ervan weglopen.

dus: kun je naar seksueel verlangen kijken zonder te vluchten, zonder te rechtvaardigen, zonder te veroordelen? genot zoekt altijd herhaling; het denken wil eraan vasthouden, en dat vasthouden wordt een vorm van gebondenheid. of het nu plezier is van seks, van roem, of gewoon van een vriendelijk woord: het denken klampt zich eraan vast, dat is de continuïteit van het ik. krishnamurti zei: de continuïteit van plezier is de continuïteit van het zelf. en dat zelf, geboren uit het denken, is altijd in conflict; het vreest verlies, vergelijkt, oordeelt. zo wordt zelfs liefde besmet door bezitterigheid en angst.

is het dan mogelijk om te leven op een manier waarin seks geen probleem is? waarin het noch verheerlijkt, noch veroordeeld wordt, maar slechts één aspect is van het menselijk bestaan, niet het middelpunt? dat is pas mogelijk als er diepe aandacht is: als het denken zich bewust is van zijn eigen bewegingen. dan begint men te zien hoe het denken het probleem creëert dat het vervolgens wil oplossen.

bijvoorbeeld: je voelt verlangen. onmiddellijk zegt het denken: dit is goed, of: dit is fout. er ontstaat verdeeldheid, het conflict wordt geboren. maar als je naar verlangen kijkt zonder inmenging, zonder de waarnemer, gebeurt er iets bijzonders. er is geen centrum om te verdedigen, geen oordeel, geen angst – slechts de energie van het waarnemen. en in dat zien verliest verlangen zijn dwingende kracht.

krishnamurti vroeg vaak: kun je kijken naar een bloem, een zonsondergang, een gezicht, zonder het verlangen om het te bezitten, om het vast te houden of te onthouden? kun je naar een ander mens kijken, niet als een object van vervulling, maar met totale genegenheid en vrijheid? dát is liefde – niet de liefde van eigendom, maar liefde zonder tegenpool. en in zo’n liefde kan er tederheid zijn, aanraking, zelfs seksualiteit, maar niet de last van verwachtingen, niet de honger om via een ander iets te worden.

daarom bood krishnamurti nooit technieken aan. geen enkele methode kan deze helderheid brengen. het vraagt om je eigen diepe observatie – niet morgen, maar nu; niet via een leraar, maar via je eigen stilte. in die stilte begin je niet alleen seks, maar de hele beweging van verlangen te begrijpen: hoe het het ik vormt, hoe het de geest bindt, hoe het volkomen doorzien kan worden.

als dat gebeurt, is er vrijheid – niet vrijheid van seks, maar van de verwarring eromheen. dan is seks niet langer heilig of profaan; het is gewoon, als de wind, als de regen, als het zonlicht. het stroomt door het leven zonder een probleem te worden, en in die eenvoud schuilt schoonheid.
